\runninghead{Scriptural Date and Location}
\section{Scriptural Date and Location}
\label{sec:date-location}

{\small Each passage the Mass reads or sings from Scripture, in canonical
order, with what is dated, where it belongs, and the date the Scripture
chronology corpus gives for it. Rows that share one answer name every element
that uses it. The chronology record also carries the Tract printed for votive
Masses after Septuagesima, which is never said on this feast; its rows are
marked so. The three orations and the Preface carry no biblical date. The
notes beneath give each passage's setting, attribution and dating evidence.\par}
\medskip

\begin{chronologyorientationtable}
\chronologyorientationrow{row-6c43dd5e0531017f}{Psalms of a king and of the Lord as King. Ps 71 is titled ``on Solomon'' and closes with the praises of David (71:20); Ps 2, where the king is set on Sion, is ascribed to David in Acts 4:25; Ps 28, which ends with the Lord sitting as King for ever, is titled ``for David, at the finishing of the tabernacle'' (28:1). No occasion is dated.}
\addlinespace
\chronologyorientationrow{row-7fbe70097919820b}{The psalms of this Mass, under the titles above. Ps 88, titled ``for Ethan the Ezrahite'' (88:1), is the psalm of the Tract, printed for votive Masses after Septuagesima, never said on this feast.}
\addlinespace
\chronologyorientationrow{row-3bd03395c37c5ada}{The Tract, printed for votive Masses after Septuagesima, never said on this feast. The setting dated is the psalm's lament that God has rejected his anointed and cast his throne down to the ground (Ps 88:39--45).}
\midrule
\chronologyorientationrow{row-94f74a8fe6bde9b6}{Daniel's night vision ``in the first year of Baltasar, king of Babylon'' (Dan 7:1), among the exiles of Juda in Babylon.}
\addlinespace
\chronologyorientationrow{row-98256791fb1a09b1}{The book that presents Daniel, a youth of Juda carried to Babylon (Dan 1:1--6). The same verse opens the paschal Alleluia, printed for Paschaltide, never said on this feast; its second verse, Apoc 19:16, has no row. A later dating is noted below.}
\midrule
\chronologyorientationrow{row-bb05ffe36adb121c}{Jesus questioned by Pilate inside the governor's hall, the praetorium, at Jerusalem, on the morning before the Pasch (Jn 18:28, 33).}
\addlinespace
\chronologyorientationrow{row-b5624d4c980f37b6}{John writes his Gospel, by early tradition at Ephesus (Irenaeus, as Fonck reports him).}
\midrule
\chronologyorientationrow{row-e679cfff41e9473d}{Paul writes in bonds (Col 4:3, 18) to the saints at Colossa (1:2); the common view places the letters of the captivity at Rome.}
\midrule
\chronologyorientationrow{row-1dc5d7a61ec35b95}{John writes the visions he saw ``in the island which is called Patmos'' (Apoc 1:9).}
\end{chronologyorientationtable}

\subsection{Psalms 2, 28 (29) and 71 (72): the Introit verse, Gradual, Offertory and Communion}
\label{sec:date-psalms}

The four chants drawn from the psalms share two rows. The ``Historical
setting'' row records that no occasion is dated for any of them. Their titles
name a king and a moment without placing either in time: Psalm~71 is \latin{in
Salomonem}, ``on Solomon,'' and its last verse closes ``the praises of David,
the son of Jesse''; Psalm~28 is ``for David, at the finishing of the
tabernacle''; Psalm~2 opens without a title, and Acts 4:25 ascribes it to
David. These attributions belong to the texts, and
the dated reign of no king is attached to any of the three psalms here. Their
witnesses also read the titles differently: for Bellarmine Psalm~71 is
literally David's prayer for Solomon, while Augustine, Cassiodorus and
Theodoret hold that it fits Solomon in no way and speaks of Christ.

The ``Composition'' row gives the bound that the introduction to the Psalms in
the \work{New American Bible, Revised Edition} sets for the whole Psalter, in
the same sentence in which it holds that no individual psalm can be dated
securely. It dates the book of Psalms as a whole and is not a date for any of
these psalms. The same row names the Tract's Psalm~88, titled for Ethan the
Ezrahite, because the chronology record carries the Tract; the Tract is
printed for votive Masses after Septuagesima and is never said on this feast.

The Tract's own ``Historical setting'' row dates the destruction of Jerusalem
under Sedecias, the setting the record gives to the psalm's lament for the
anointed rejected and the throne cast down (Ps 88:39--45). Its preferred
figure is the one Haydock's Douay commentary prints in a chronological
paragraph attributed to Ussher; the alternatives are those of the
\work{Catholic Encyclopedia}'s articles ``Captivities of the Israelites''
(Reid) and ``Temple of Jerusalem'' (Meistermann), ``Third and Fourth Books of
Kings'' (Schets), and ``Chronology of the Kings'' (Sloet, reporting Petavius).
They date a setting of a chant that this feast does not sing.

\subsection{Daniel 7:14: the Alleluia}
\label{sec:date-daniel}

Daniel dates the vision itself to the first year of Baltasar, king of Babylon
(7:1), but the chronology record holds no row for that year. The ``Historical
setting'' row gives instead the seventy years of servitude to Babylon that
Jeremias foretold (Jer 25:11--12), as a duration and not as a year. The book
presents its author as Daniel, one of the youths of Juda carried to Babylon
and taught the learning of the Chaldeans (Dan 1:1--6).

The ``Composition'' row gives the date that Francis Gigot's ``Book of Daniel''
in the \work{Catholic Encyclopedia} (vol.~4) calls the time-honoured position.
Gigot sets beside it what he calls a comparatively recent theory, widely
accepted by contemporary scholars, which dates the book in the time of
Antiochus IV Epiphanes, and he concludes that Catholic scholars generally
abide by the traditional view, although they are not bound to it by any
decision of the Church. The traditional answer the row displays excludes the
Antiochene date, and the critical answer the record holds offers no dated
alternative for the book.

The paschal Alleluia printed for Paschaltide opens with this same verse of
Daniel and adds Apoc 19:16, ``KING OF KINGS AND LORD OF LORDS.'' It is never
said on this feast, and the chronology record has no element for it. The
Apocalypse composition dates that the Introit's verses receive below are not
extended by the record to Apoc 19:16.

\subsection{John 18:33--37: the Gospel}
\label{sec:date-john}

The ``Event'' row has no date. The chronology record dates the crucifixion and
binds that date to the passion narrative from Jn 19:18, but no event or binding
reaches the questioning before Pilate in Jn 18:33--37, and none is supplied
here. John places the scene in the governor's hall at Jerusalem, early in the
morning, on the day before the Pasch (18:28). The composition dates shown are
those of Alfred Durand's ``The New Testament'' (vol.~14) and Leopold Fonck's
``Gospel of St.~John'' (vol.~8) in the \work{Catholic Encyclopedia}; both are
marked disputed. Fonck reports St Irenaeus (\work{Adversus Haereses} III.1.2)
for the tradition that John wrote his Gospel at Ephesus. The dates are dates of
the book, not of the trial.

\subsection{Colossians 1:12--20: the Epistle}
\label{sec:date-colossians}

The two composition dates are those of Ferdinand Prat's ``St.~Paul''
(vol.~11), which gives one year to the letters of the captivity, and Cornelius
Aherne's ``Epistle to the Colossians'' (vol.~4), which gives a range; both
are marked disputed. Aherne reports the common view that the captivity
letters were written at Rome. The letter itself says that Paul is bound (4:3,
18) and is addressed to the saints at Colossa (1:2): the writer's place and the
recipients' place are not the same. The passage narrates no event.

\subsection{Apocalypse 1:6 and 5:12: the Introit}
\label{sec:date-apocalypse}

The two composition dates are those of Christian Van den Biesen's
``Apocalypse'' (vol.~1), which places the book in the reign of Domitian, and of
Durand's ``The New Testament''; both are marked disputed. John writes that he
was ``in the island which is called Patmos'' when he saw the visions (1:9),
in the same chapter as the Introit's verse~6. The throne vision of chapters
4--5 is a vision and carries no date of its own.
